%!TeX root = ../main.tex

\subsection{Análisis de estabilidad del sistema a través del criterio de Lyapunov}

El criterio de Lyapunov implica la existencia de una función $\mathbf{V}(x)$ considerada
función candidata, dicha función definida en \eqref{eq:DefiniciónLyapunovV} requiere de
una matriz $P$ que sea positiva definida.
Dicha función $\mathbf{V}(x)$ deberá ser definida positiva y su derivada temporal
$\mathbf{\dot{V}}(x)$ deberá definida negativa o semidefinida negativa.

\begin{equation}
  \begin{aligned}
    \mathbf{V}(x)       &= \mathbf{x}^TP\mathbf{x},\qquad P=P^T > 0  \\
    \mathbf{\dot{V}}(x) &= \mathbf{x}^T \left(\mathbf{A} P + P \mathbf{A} \right)\mathbf{x}
    \label{eq:DefiniciónLyapunovV}
  \end{aligned}
\end{equation}

\subsubsection{Función candidata de Lyapunov}

Para simplificar el análisis, se considera una matriz $P$ diagonal con elementos
positivos \eqref{eq:PDiagonal}.

\begin{equation}
    P = \begin{bmatrix}
        \alpha & 0 & 0 \\
        0 & \beta & 0 \\
        0 & 0 & \gamma
    \end{bmatrix}, \quad \alpha, \beta, \gamma > 0
    \label{eq:PDiagonal}
\end{equation}

\subsubsection{Desarrollo de la función candidata}
Sustituyendo la matriz diagonal $P$ de \eqref{eq:PDiagonal} y el vector de estado $\mathbf{x} = [i, \theta, \omega]^T$ en la definición general \eqref{eq:DefiniciónLyapunovV}, se obtiene la función candidata concreta.
\begin{equation*}
    \begin{aligned}
        \mathbf{V}(\mathbf{x}) &= \mathbf{x}^T P \mathbf{x}\\
        \mathbf{V}(\mathbf{x}) &=
        \begin{bmatrix}
            i & \theta & \omega
        \end{bmatrix}
        \begin{bmatrix}
            \alpha & 0 & 0 \\
            0 & \beta & 0 \\
            0 & 0 & \gamma
        \end{bmatrix}
        \begin{bmatrix}
            i \\ \theta \\ \omega
        \end{bmatrix}
    \end{aligned}
\end{equation*}
\begin{equation}
        \mathbf{V}(\mathbf{x}) =
        \alpha i^2 + \beta \theta^2 + \gamma \omega^2
        \label{eq:LyapunovCandidata}
\end{equation}

Sustituyendo la matriz $\mathbf{A}$ y la matriz $P$ diagonal en la expresión $\mathbf{A}^T P + P \mathbf{A}$ se obtiene la expresión \eqref{eq:ATPPA}.

\begin{equation}
    \mathbf{A}^T P + P \mathbf{A} = 
    \begin{bmatrix} 
        -2\frac{R}{L}\alpha & 0 & \frac{K_t}{J}\gamma - \frac{K_e}{L}\alpha \\
        0 & 0 & \beta \\
        \frac{K_t}{J}\gamma - \frac{K_e}{L}\alpha & \beta & -2\frac{B}{J}\gamma
    \end{bmatrix}
    \label{eq:ATPPA}
\end{equation}

\subsubsection{Condición para simplificar la matriz}

Para eliminar los términos cruzados (1,3) y (3,1) se impone la condición
\eqref{eq:CondicionCruzados}.

\begin{equation}
    \frac{K_t}{J}\gamma - \frac{K_e}{L}\alpha = 0
    \label{eq:CondicionCruzados}
\end{equation}

Asumiendo que $K_e = K_t$ en unidades consistentes\footnote{Esto se puede considerar
porque en un motor DC ideal, sin pérdidas, la potencia eléctrica que entra es igual
a la potencia mecánica que sale. Es decir, el voltaje contraelectromotriz por la
corriente es igual al par por la velocidad.}, se obtiene la relación
\eqref{eq:GammaAlpha}.

\begin{equation}
    \gamma = \frac{J}{L} \alpha
    \label{eq:GammaAlpha}
\end{equation}

Sustituyendo \eqref{eq:GammaAlpha} en \eqref{eq:ATPPA}, la matriz se simplifica
a \eqref{eq:ATPPASimplificada}.

\begin{equation}
    \mathbf{A}^T P + P \mathbf{A} = 
    \begin{bmatrix} 
        -2\frac{R}{L}\alpha & 0 & 0 \\
        0 & 0 & \beta \\
        0 & \beta & -2\frac{B}{L}\alpha
    \end{bmatrix}
    \label{eq:ATPPASimplificada}
\end{equation}

\subsubsection{Derivada de la función de Lyapunov}

La derivada temporal de la función de Lyapunov resulta en la expresión
\eqref{eq:DotVFinal}.

\begin{equation}
    \mathbf{\dot{V}}(\mathbf{x}) = -2\alpha\left( \frac{R}{L} i^2 + \frac{B}{L} \omega^2 \right) + 2\beta \theta \omega
    \label{eq:DotVFinal}
\end{equation}

\subsubsection{Verificación de las condiciones de Lyapunov}

La función $\mathbf{V}(\mathbf{x}) = \alpha i^2 + \beta \theta^2 + \gamma \omega^2$ es
definida positiva para $\alpha, \beta, \gamma > 0$. Sin embargo, la derivada
$\mathbf{\dot{V}}(\mathbf{x})$ no es definida negativa para todos los estados debido al
término $2\beta \theta \omega$. Si se considera $\beta = 0$, la derivada se
reduce a \eqref{eq:DotVSimplificada}.

\begin{equation}
    \mathbf{\dot{V}}(\mathbf{x}) = -2\alpha\left( \frac{R}{L} i^2 + \frac{B}{L} \omega^2 \right) \leq 0
    \label{eq:DotVSimplificada}
\end{equation}

La función candidata de Lyapunov en \eqref{eq:LyapunovCandidata} es definida positiva
y su derivada en \eqref{eq:DotVSimplificada} es semidefinida negativa, lo que
satisface las condiciones para estabilidad en el sentido de Lyapunov. Esta
estabilidad no es asintótica, pues $\mathbf{\dot{V}}(\mathbf{x})$ se anula para
$i=0$, $\omega=0$ y cualquier valor de $\theta$, lo cual refleja directamente
el autovalor en cero de la matriz dinámica $\mathbf{A}$.
